% ============================================================================
% MÓDULO 2 — SIMULADO E REVISÃO INTEGRADA
% CHF662 — ANÁLISE AMBIENTAL (DCHF/UEFS)
% ============================================================================

\clearpage
\phantomsection
\addcontentsline{toc}{chapter}{Simulado e Revisão Integrada — Módulo 2}
\chapter*{Simulado e Revisão Integrada — Módulo 2}
\markboth{SIMULADO — MÓDULO 2}{SIMULADO — MÓDULO 2}

\begin{boxConceito}[Instruções de Aplicação e Autodiagnóstico (Semanas 5 e 6)]
  Este simulado integra integralmente os conteúdos dos \textbf{Capítulos 5 e 6} (Avaliação de Impacto Ambiental, Matriz de Leopold, Hierarquia de Mitigação, PBA, Sistemas Ambientais Abertos, Escalas, Falácia de Escala e Ordenamento Territorial/ZEE), preparando você simultaneamente para o \textbf{Estudo Orientado EO4}, para a \textbf{Avaliação Teórica da Semana 10} e para a etapa metodológica do \textbf{Trabalho Prático Final}.
  \begin{itemize}[leftmargin=1.5em, itemsep=2pt]
    \item \textbf{Tempo sugerido:} $90$ minutos ($40$ min para as 8 questões objetivas + $50$ min para as 2 questões dissertativas).
    \item \textbf{Regra de ouro:} resolva todas as questões sem consultar o gabarito; somente depois faça a correção comentada e preencha a \textbf{Tabela de Diagnóstico de Revisão} ao final.
  \end{itemize}
\end{boxConceito}

\section*{Parte I — Questões Objetivas (Estilo ENADE e Concursos Ambientais)}

\begin{enumerate}[leftmargin=1.8em, itemsep=6pt]
  \item \textbf{[Autoral — Estilo ENADE / AIA e Etapas]} Um empreendedor solicitou licença ambiental para construir uma pequena central hidrelétrica em um rio intermitente do interior baiano. Antes de iniciar o levantamento de campo, o órgão ambiental verificou o porte e a sensibilidade da área, definiu que seria exigido EIA/RIMA e emitiu o documento contendo as diretrizes mínimas de estudo. Essas duas etapas iniciais correspondem, respectivamente, a:
  \begin{alternativas}
    \item Prognóstico e Plano Básico Ambiental (PBA).
    \item Triagem (\textit{screening}) e Definição de Escopo (\textit{scoping} / Termo de Referência).
    \item Monitoramento pós-implantação e Compensação Ambiental.
    \item Diagnóstico da linha de base e Auditoria de Licença de Operação.
    \item Sobreposição cartográfica e Rede de Interação.
  \end{alternativas}
  \espacoResolucao[1.1cm]

  \item \textbf{[Autoral — Estilo Concurso Analista Ambiental / Matriz de Leopold]} Em uma Matriz de Leopold Ponderada ($P_{ij} = s_{ij} \cdot M_{ij} \cdot I_{ij}$), a ação ``escavação de jazida de empréstimo'' apresentou três interações: estabilidade de encostas $(-, M=9, I=8)$, qualidade da água por sedimentos $(-, M=5, I=6)$ e renda de prestadores locais $(+, M=6, I=5)$. Assinale a alternativa correta sobre a pontuação e a interpretação técnica dessa coluna:
  \begin{alternativas}
    \item $\Sigma^{-} = -102$, $\Sigma^{+} = +30$, total $T = -72$; a renda gerada não elimina o risco geotécnico de $-72$ em estabilidade de encostas.
    \item $\Sigma^{-} = -72$, $\Sigma^{+} = +60$, total $T = -12$; a obra dispensa mitigação por saldo quase neutro.
    \item $\Sigma^{-} = -102$, $\Sigma^{+} = +30$, total $T = -132$; os valores positivos devem ser subtraídos da carga negativa.
    \item $\Sigma^{-} = -28$, $\Sigma^{+} = +11$, total $T = -17$; magnitude e importância devem ser somadas em cada célula.
    \item $\Sigma^{-} = -72$, $\Sigma^{+} = +30$, total $T = -42$; o impacto na água é anulado pelo impacto positivo em renda.
  \end{alternativas}
  \espacoResolucao[1.2cm]

  \item \textbf{[Autoral — Estilo Concurso Órgão Ambiental / Áreas de Influência e Métodos]} Uma equipe técnica precisa avaliar um projeto de linha de transmissão que atravessa três municípios do Semiárido baiano. Para escolher o traçado com menor conflito sobre APPs e remanescentes de Caatinga e, simultaneamente, mapear efeitos indiretos sobre a erosão e os açudes a jusante, a combinação metodológica mais consistente consiste em:
  \begin{alternativas}
    \item Delimitar a AID como círculo fixo de $1\,\text{km}$ e aplicar exclusivamente uma listagem de controle simples.
    \item Integrar sobreposição de cartas em SIG (para alternativas locacionais) com redes de interação (para cadeias de impactos diretos e indiretos), delimitando a AID física pelas microbacias interceptadas.
    \item Dispensar a linha de base biótica e aplicar apenas a compensação financeira prevista no SNUC.
    \item Utilizar apenas a AII municipal para projetar a fundação de cada torre, desconsiderando a ADA.
    \item Substituir o prognóstico comparativo pelo cenário estático anterior à obra.
  \end{alternativas}
  \espacoResolucao[1.1cm]

  \item \textbf{[Autoral — Estilo Concurso IBAMA / Hierarquia de Mitigação]} O projeto original de ampliação de um distrito industrial previa suprimir $16{,}0\,\text{ha}$ de vegetação nativa. Após revisão pela Hierarquia de Mitigação: (1) o redesenho das quadras evitou $25\%$ da supressão inicial; (2) técnicas construtivas seletivas minimizaram $25\%$ da área restante; e (3) a revegetação de áreas temporárias recuperou $3{,}0\,\text{ha}$. Se o órgão licenciador exigir compensação na razão $2:1$ sobre o impacto permanente residual, a área total a ser compensada será de:
  \begin{alternativas}
    \item $6{,}0\,\text{ha}$.
    \item $9{,}0\,\text{ha}$.
    \item $12{,}0\,\text{ha}$.
    \item $18{,}0\,\text{ha}$.
    \item $32{,}0\,\text{ha}$.
  \end{alternativas}
  \espacoResolucao[1.4cm]

  \item \textbf{[Autoral — Estilo ENADE / Sistemas Abertos e Retroalimentação]} Sobre o funcionamento de paisagens e bacias hidrográficas como sistemas ambientais abertos (Drew, 1989; Castro; Lemos, 2016), analise as afirmativas:
  \begin{enumerate}[label=\Roman*., leftmargin=2.0em, itemsep=2pt]
    \item Quando as entradas superam as saídas ($\Delta S > 0$) no balanço de sedimentos de um reservatório, ocorre acúmulo interno de sedimentos (assoreamento), reduzindo o volume útil de água.
    \item A retroalimentação negativa atua amortecendo desvios em relação ao equilíbrio dinâmico, enquanto a retroalimentação positiva amplifica perturbações e pode acelerar processos de degradação.
    \item Um limiar (\textit{threshold}) corresponde ao ponto crítico a partir do qual o sistema sofre mudança qualitativa ou abrupta em seu regime de funcionamento.
  \end{enumerate}
  É correto o que se afirma em:
  \begin{alternativas}
    \item I, apenas.
    \item II, apenas.
    \item I e II, apenas.
    \item II e III, apenas.
    \item I, II e III.
  \end{alternativas}
  \espacoResolucao[1.1cm]

  \item \textbf{[Autoral — Estilo Concurso Geógrafo / Falácia de Escala e SDR]} Em uma bacia hidrográfica de $1\,000\,\text{ha}$, a erosão bruta estimada em parcelas nas vertentes é de $20{,}0\,\text{t/(ha}\cdot\text{ano)}$, mas o monitoramento no exutório registra uma produção total de sedimentos de $5\,000\,\text{t/ano}$. Sobre esse diagnóstico, é correto afirmar que:
  \begin{alternativas}
    \item A Razão de Entrega de Sedimentos é $\text{SDR} = 0{,}25$, indicando que $15\,000\,\text{t/ano}$ ($75\%$) ficam retidas nas vertentes, sopés e planícies; assumir que todo o solo erodido chega ao exutório configura falácia de escala.
    \item A Razão de Entrega de Sedimentos é $\text{SDR} = 0{,}75$, pois $75\%$ dos sedimentos chegam ao exutório.
    \item A erosão bruta total da bacia é de $5\,000\,\text{t/ano}$, havendo igualdade entre a escala de parcela e a escala de bacia.
    \item O monitoramento no exutório está incorreto, porque em sistemas abertos $\text{SDR}$ é sempre maior que $1$.
    \item A retenção interna de sedimentos é de $5\,000\,\text{t/ano}$, igual à carga exportada.
  \end{alternativas}
  \espacoResolucao[1.4cm]

  \item \textbf{[Autoral — Estilo ENADE / Escala Cartográfica e Análise Integrada]} Um município pretende utilizar a base cartográfica do Zoneamento Ecológico-Econômico estadual na escala $1:250\,000$ como única fonte para delimitar, em campo, faixas de mata ciliar de $30\,\text{m}$ e ravinas de $5\,\text{m}$ de largura. Considerando o erro gráfico mínimo de $0{,}2\,\text{mm}$ no papel, essa decisão técnica é:
  \begin{alternativas}
    \item Adequada, pois na escala $1:250\,000$ a menor feição representável é de $0{,}5\,\text{m}$.
    \item Adequada para a mata ciliar e inadequada apenas para as ravinas.
    \item Inadequada para ambos os alvos, pois na escala $1:250\,000$ a menor feição representável é de $50\,\text{m}$, incorrendo em falácia de escala por uso de base regional em demanda de detalhe local.
    \item Adequada, porque a escala cartográfica não interfere na visualização de feições lineares.
    \item Inadequada apenas porque o ZEE não considera a vegetação ripária.
  \end{alternativas}
  \espacoResolucao[1.2cm]

  \item \textbf{[Autoral — Estilo Concurso Planejamento Territorial / Ordenamento e ZEE]} Assinale a alternativa que articula corretamente \textbf{planejamento ambiental}, \textbf{planejamento territorial}, \textbf{ordenamento territorial} e \textbf{ZEE}:
  \begin{alternativas}
    \item O ZEE dispensa o planejamento ambiental porque trata exclusivamente da arrecadação tributária municipal.
    \item O planejamento ambiental estabelece a leitura dos limites, fragilidades e serviços ecossistêmicos; o planejamento territorial integra essa base às estratégias socioespaciais e de infraestrutura; e o ordenamento territorial confere expressão espacial normativa por meio de instrumentos como o ZEE.
    \item O ordenamento territorial antecede qualquer diagnóstico geoambiental e independe da escala adotada.
    \item O planejamento territorial restringe-se às áreas urbanas, enquanto o ordenamento territorial aplica-se apenas a Unidades de Conservação de proteção integral.
    \item Planejamento ambiental e ordenamento territorial são instrumentos exclusivos da iniciativa privada no âmbito do PBA.
  \end{alternativas}
  \espacoResolucao[1.1cm]
\end{enumerate}

\section*{Parte II — Questões Dissertativas (Formato da Avaliação Teórica CHF662)}

\begin{enumerate}[leftmargin=1.8em, itemsep=6pt, resume]
  \item \textbf{[Questão Dissertativa 1 — AIA, Matriz Ponderada, Hierarquia de Mitigação e PBA (Valor de referência: $1{,}5$ ponto)]}
  Uma mineradora de brita planeja instalar uma pedreira próxima a uma microbacia do Semiárido baiano. Na avaliação preliminar, cruzaram-se duas ações ($A_1$: decapeamento do solo e corte de vegetação; $A_2$: implantação de bacias de contenção e operação regular) com três fatores ($F_1$: solo e processos erosivos; $F_2$: vegetação e APP; $F_3$: empregos e serviços locais), obtendo-se as pontuações $(s, M, I)$:
  \begin{itemize}[leftmargin=1.5em, itemsep=1.5pt]
    \item $F_1$: $A_1=(-, 8, 9)$; $A_2=(-, 3, 5)$;
    \item $F_2$: $A_1=(-, 7, 9)$; $A_2=(-, 2, 6)$;
    \item $F_3$: $A_1=(+, 4, 5)$; $A_2=(+, 6, 7)$.
  \end{itemize}
  \begin{enumerate}[label=(\alph*), leftmargin=1.8em, itemsep=2pt]
    \item Calcule todas as células $P_{ij}$, os totais por linha ($F_i$), os totais por coluna ($T_j$), as cargas $\Sigma^{-}$ e $\Sigma^{+}$ e o balanço $B$. Justifique por que o saldo positivo em $F_3$ não anula a criticidade de $F_1$ e $F_2$.
    \item Aplique os quatro degraus da \textbf{Hierarquia de Mitigação} propondo uma ação concreta para cada degrau no caso da pedreira.
    \item Indique dois programas obrigatórios do \textbf{Plano Básico Ambiental (PBA)} e os respectivos indicadores mensuráveis que devem ser comparados à linha de base no monitoramento pós-implantação.
  \end{enumerate}
  \espacoResolucao[4.0cm]

  \item \textbf{[Questão Dissertativa 2 — Sistemas Abertos, Limiares, Escalas e Ordenamento Territorial (Valor de referência: $1{,}5$ ponto)]}
  Em uma microbacia semiárida degradada por sobrepastejo, o monitoramento anual registrou precipitação $P = 720\,\text{mm/ano}$, evapotranspiração real $ET = 480\,\text{mm/ano}$, escoamento superficial $Q_s = 260\,\text{mm/ano}$ e escoamento de base $Q_b = 20\,\text{mm/ano}$. Durante um evento convectivo de verão, caiu uma chuva de intensidade $I = 70\,\text{mm/h}$ sobre um solo compactado cuja capacidade de infiltração era $f_c = 25\,\text{mm/h}$.
  \begin{enumerate}[label=(\alph*), leftmargin=1.8em, itemsep=2pt]
    \item Calcule a variação do estoque hídrico anual ($\Delta S$) e o excesso de precipitação convertido em enxurrada ($R$) durante o evento chuvoso (em $\text{mm/h}$ e em percentual da chuva). Explique o papel do limiar de infiltração e da retroalimentação positiva na degradação observada.
    \item Explique como a escala espacial e a escala temporal diferem entre o diagnóstico do evento chuvoso na vertente e o planejamento da sub-bacia regional, indicando o risco de falácia de escala.
    \item Diferencie \textbf{planejamento ambiental}, \textbf{planejamento territorial} e \textbf{ordenamento territorial}, explicando como o \textbf{Zoneamento Ecológico-Econômico (ZEE)} deve enquadrar essa microbacia degradada.
  \end{enumerate}
  \espacoResolucao[4.0cm]
\end{enumerate}

% ============================================================================
% GABARITO COMENTADO, ESPELHO DISSERTATIVO E DIAGNÓSTICO DE REVISÃO
% ============================================================================
\clearpage
\section*{Gabarito Comentado — Simulado do Módulo 2}

\begin{center}
\renewcommand{\arraystretch}{1.2}
\small
\begin{tabularx}{\textwidth}{@{}c c c c c c c c@{}}
  \toprule
  \rowcolor{uefsNavy}
  \textcolor{white}{\bfseries Q1} & \textcolor{white}{\bfseries Q2} & \textcolor{white}{\bfseries Q3} & \textcolor{white}{\bfseries Q4} & \textcolor{white}{\bfseries Q5} & \textcolor{white}{\bfseries Q6} & \textcolor{white}{\bfseries Q7} & \textcolor{white}{\bfseries Q8} \\
  \midrule
  \textbf{B} & \textbf{A} & \textbf{B} & \textbf{C} & \textbf{E} & \textbf{A} & \textbf{C} & \textbf{B} \\
  \bottomrule
\end{tabularx}
\end{center}

\begin{boxGabarito}[Comentário Detalhado das Questões Objetivas (Q1 a Q8)]
  \begin{itemize}[leftmargin=1.5em, itemsep=3.5pt]
    \item \textbf{Q1 — (B):} decidir pela exigência de EIA/RIMA segundo porte e sensibilidade é \textbf{triagem (\textit{screening})}; emitir o Termo de Referência fixando diretrizes e profundidade é \textbf{definição de escopo (\textit{scoping})}.
    \item \textbf{Q2 — (A):} células: $-(9 \times 8) = -72$, $-(5 \times 6) = -30$ e $+(6 \times 5) = +30$. Assim, $\Sigma^{-} = -72 - 30 = -102$, $\Sigma^{+} = +30$ e $T = -102 + 30 = -72$. Benefícios de renda são incomensuráveis com risco geotécnico de encosta.
    \item \textbf{Q3 — (B):} SIG/sobreposição orienta alternativas locacionais; redes de interação evidenciam impactos indiretos em cascata; a AID física deve respeitar as microbacias interceptadas, não círculos arbitrários.
    \item \textbf{Q4 — (C):} passo a passo na Hierarquia de Mitigação: inicial $= 16{,}0\,\text{ha}$; 1º Evitar ($25\%$): $16{,}0 - 4{,}0 = 12{,}0\,\text{ha}$; 2º Minimizar ($25\%$ de $12{,}0$): $12{,}0 - 3{,}0 = 9{,}0\,\text{ha}$; 3º Recuperar ($3{,}0\,\text{ha}$ temporários): residual permanente $= 9{,}0 - 3{,}0 = 6{,}0\,\text{ha}$; 4º Compensar ($2:1$): $2 \times 6{,}0 = \mathbf{12{,}0\,\text{ha}}$.
    \item \textbf{Q5 — (E):} todas as três afirmativas (I, II e III) expressam corretamente os princípios de balanço de massa, retroalimentação negativa/positiva e limiares em sistemas ambientais abertos.
    \item \textbf{Q6 — (A):} erosão bruta total $= 20{,}0 \times 1\,000 = 20\,000\,\text{t/ano}$. Como chegam $5\,000\,\text{t/ano}$ ao exutório, $\text{SDR} = 5\,000 / 20\,000 = \mathbf{0{,}25}$ ($25\%$), restando $20\,000 - 5\,000 = \mathbf{15\,000\,\text{t/ano}}$ ($75\%$) retidos internamente na bacia.
    \item \textbf{Q7 — (C):} $d_{\min} = 0{,}0002 \times 250\,000 = \mathbf{50\,\text{m}}$. Tanto a mata ciliar ($30\,\text{m}$) quanto as ravinas ($5\,\text{m}$) estão abaixo da resolução gráfica da carta, caracterizando falácia de escala descendente.
    \item \textbf{Q8 — (B):} apresenta com rigor a articulação complementar entre planejamento ambiental, planejamento territorial e ordenamento territorial/ZEE (Arruda, 2013; IBGE, 2018).
  \end{itemize}
\end{boxGabarito}

\section*{Espelho de Correção das Questões Dissertativas (Q9 e Q10)}

\begin{boxGabarito}[Espelho Dissertativo — Questão 9 ($1{,}5$ ponto)]
  \begin{itemize}[leftmargin=1.5em, itemsep=3pt]
    \item \textbf{Item (a) — $0{,}6$ ponto:}
    Células: $P_{11} = -(8 \times 9) = -72$, $P_{12} = -(3 \times 5) = -15 \implies F_1 = -87$; $P_{21} = -(7 \times 9) = -63$, $P_{22} = -(2 \times 6) = -12 \implies F_2 = -75$; $P_{31} = +(4 \times 5) = +20$, $P_{32} = +(6 \times 7) = +42 \implies F_3 = +62$.
    Colunas: $T_1 = -72 - 63 + 20 = -115$; $T_2 = -15 - 12 + 42 = +15$.
    Cargas e balanço: $\Sigma^{-} = -72 - 15 - 63 - 12 = -162$; $\Sigma^{+} = +20 + 42 = +62$; $B = -162 + 62 = -100$ (conferência: $-87 - 75 + 62 = -100$ e $-115 + 15 = -100$).
    \emph{Justificativa técnica:} os $+62$ pontos de emprego pertencem à esfera econômica e não restauram fisicamente o solo erodido ($-87$) nem a vegetação/APP ($-75$), exigindo mitigação dirigida às células críticas ($P_{11} = -72$ e $P_{21} = -63$).
    \item \textbf{Item (b) — $0{,}5$ ponto:}
    \textbf{1º Evitar:} redefinir a poligonal de lavra e acessos para manter intacta a APP e trechos mais íngremes; \textbf{2º Minimizar:} implantar drenagem perimetral, bacias de decantação, aspersão de poeira e decapeamento em fases; \textbf{3º Recuperar:} estocar o solo superficial (\textit{topsoil}) e executar PRAD com bioengenharia e espécies nativas da Caatinga nas bancadas exauridas; \textbf{4º Compensar:} reposição florestal ou apoio a UC proporcional apenas à área de perda permanente residual.
    \item \textbf{Item (c) — $0{,}4$ ponto:}
    Exemplos aceitos: (1) \emph{Programa de Monitoramento da Qualidade da Água e Sedimentos} (indicadores: turbidez em NTU e sólidos suspensos totais em $\text{mg/L}$ comparados à linha de base); (2) \emph{Programa de Controle Erosivo e PRAD} (indicadores: taxa de sobrevivência de mudas, cobertura vegetal em $\%$ e monitoramento de sulcos/pinos de erosão).
  \end{itemize}
\end{boxGabarito}

\begin{boxGabarito}[Espelho Dissertativo — Questão 10 ($1{,}5$ ponto)]
  \begin{itemize}[leftmargin=1.5em, itemsep=3pt]
    \item \textbf{Item (a) — $0{,}6$ ponto:}
    Balanço hídrico anual: $S_{\text{out}} = 480 + 260 + 20 = 760\,\text{mm/ano} \implies \Delta S = 720 - 760 = \mathbf{-40\,\text{mm/ano}}$ (depleção de estoque).
    Evento chuvoso: como $I = 70\,\text{mm/h} > f_c = 25\,\text{mm/h}$, rompe-se o limiar extrínseco de infiltração, gerando enxurrada $R = 70 - 25 = \mathbf{45\,\text{mm/h}}$, o que equivale a $\frac{45}{70} \times 100\% \approx \mathbf{64{,}29\%}$ da chuva.
    \emph{Retroalimentação positiva:} o excesso de enxurrada arrasta solo superficial, matéria orgânica e sementes, agravando o selamento, reduzindo ainda mais $f_c$ e amplificando o ciclo de degradação e baixa recarga ($Q_b = 20\,\text{mm/ano}$).
    \item \textbf{Item (b) — $0{,}45$ ponto:}
    O evento chuvoso na vertente opera na microescala/escala de vertente e na escala temporal de minutos a horas (evento), governado por infiltração e erosão concentrada. O planejamento da sub-bacia regional opera na mesoescala/macroescala e em escalas de anos a décadas, integrando retenção de sedimentos ($\text{SDR} < 1$), uso da terra e gestão hídrica. Extrapolar linearmente o escoamento ou a erosão da vertente para toda a sub-bacia (ou usar mapas regionais grosseiros para projetar obras locais) incorre em \textbf{falácia de escala}.
    \item \textbf{Item (c) — $0{,}45$ ponto:}
    \emph{Planejamento ambiental:} diagnostica fragilidades, estoques hídricos e degradação; \emph{planejamento territorial:} articula esses limites às atividades pecuárias, estradas e demandas comunitárias; \emph{ordenamento territorial:} normatiza espacialmente o uso da terra. No \textbf{ZEE}, essa microbacia degradada deve ser enquadrada prioritariamente como \textbf{Zona de Recuperação Ambiental} (ou manejo restrito pós-reabilitação), recebendo ações de exclusão temporária do sobrepastejo, conservação de solo/água e recomposição vegetal.
  \end{itemize}
\end{boxGabarito}

\section*{Tabela de Diagnóstico de Revisão — Módulo 2}

Marque os itens em que você errou ou teve insegurança e execute imediatamente a revisão direcionada indicada na última coluna antes de avançar ao Módulo 3:

\begin{center}
\renewcommand{\arraystretch}{1.22}
\small
\begin{tabularx}{\textwidth}{@{}c c p{4.2cm} >{\raggedright\arraybackslash}X@{}}
  \toprule
  \rowcolor{uefsNavy}
  \textcolor{white}{\bfseries Questão} & \textcolor{white}{\bfseries Cap.} & \textcolor{white}{\bfseries Conteúdo Avaliado} & \textcolor{white}{\bfseries Ação Corretiva de Revisão} \\
  \midrule
  Q1, Q3 & Cap. 5 & Etapas da AIA, ADA/AID/AII e métodos & Reler Seção 5.2 (Figura~\ref{fig:aia_mitigacao}), Exemplo 5.1 e Erro Comum 4. \\
  \rowcolor{bgGray}
  Q2, Q9(a) & Cap. 5 & Matriz de Leopold Ponderada e incomensurabilidade & Refazer o quadro de equações do Cap. 5, o Exemplo 5.2 e o Erro Comum 2. \\
  Q4, Q9(b--c) & Cap. 5 & Hierarquia de Mitigação, PBA e monitoramento & Revisar Painel B da Figura~\ref{fig:aia_mitigacao}, Exemplos 5.4 e 5.5 e Erro Comum 3. \\
  \rowcolor{bgGray}
  Q5, Q10(a) & Cap. 6 & Sistemas abertos, balanço, limiares e retroalimentação & Reler Seção 6.2 (Figura~\ref{fig:sistema_escalas}), Exemplos 6.1 e 6.3 e Erro Comum 1. \\
  Q6, Q7, Q10(b) & Cap. 6 & Escalas, SDR, resolução cartográfica e falácia de escala & Refazer o quadro de fórmulas do Cap. 6, Exemplos 6.2 e 6.4 e Erro Comum 2. \\
  \rowcolor{bgGray}
  Q8, Q10(c) & Cap. 6 & Planejamento ambiental, territorial, ordenamento e ZEE & Revisar a tabela comparativa da Seção 6.4, Seção 6.5, Exemplo 6.5 e Erro Comum 4. \\
  \bottomrule
\end{tabularx}
\end{center}
